% =======================================================================================
% Chapter 5 -- Conclusions
% SECTION-H(14): a detailed evaluation of results and an overall assessment of the outcomes.
% =======================================================================================
\chapter{Conclusions}
\label{ch:conclusions}

\section{Overall Assessment}
\label{sec:assessment}

This study set out to build a bilingual English--Spanish voice transcription and synthesis system
that preserves the speaker's identity, and to establish how much of the speaker's prosody it
carries. Both were achieved, and the second produced a result that is more informative than the
first.

The system works and is reliable: 868 of 870 held-out utterances completed, the two failures having
a single identified cause, and the synthesised speech is easy to re-recognise (corpus
intelligibility WER 0.081 and 0.087). Speaker identity is preserved well --- 94.5\,\% and 86.1\,\%
of the human cross-language ceiling. Prosody is not: within a speaker, pitch level transfers at
$0.31$ against a human $0.85$, and pitch range essentially not at all.

The strength of the study lies less in the system than in what the corpus permitted. Because one
bilingual speaker performs both sides of every pair, each measurement has a human answer beside it.
That turns two uninterpretable numbers into calibrated ones and separates two properties that are
routinely conflated: a pipeline can preserve \emph{who is speaking} very well while preserving
\emph{how they spoke} very poorly, and only a same-speaker reference makes that visible.

\section{Conclusions Against the Objectives}
\label{sec:objectives_met}

\textbf{Objective 1 --- bidirectional voice-preserving translation framework: achieved.} A cascade
of Whisper, MarianMT and XTTS-v2 was implemented and evaluated in both directions on a
speaker-disjoint split, with deterministic per-utterance seeding and every failure retained in the
results. Identity preservation was measured against a human cross-language anchor rather than
reported as an unanchored cosine, and reaches 94.5\,\% (\ento{}) and 86.1\,\% (\esto{}) of that
ceiling. The asymmetry between directions is a property of the synthesiser, since the two ceilings
are identical to three decimal places.

\textbf{Objective 2 --- prosody analysis, transfer measurement and prediction: achieved, with a
negative result.} The human cross-lingual prosodic relationship was characterised over 2{,}304
same-speaker pairs, decomposed into within- and between-speaker components, and used as the
reference for the system. The transfer gap was quantified feature by feature. A supervised mapping
from source to target prosody improved on a copy-plus-offset baseline only for pitch range, by
5.1\,\% and 3.9\,\% of mean absolute error, surviving family-wise correction in both directions;
pitch level and timing were not predictable beyond the baseline. The learned coefficient on source
pitch range is negative, which supplies the mechanism: utterance-level variation regresses towards
the speaker's typical range rather than transferring one for one.

\textbf{Objective 3 --- integration with video: achieved, with a stated limit.} A web application
extracts audio, transcribes with timestamps, translates and synthesises each segment in the
original speaker's voice, retimes each segment to its slot with WORLD analysis and resynthesis, and
remuxes without re-encoding the picture. On five real jobs the dubbed track lands within 7\,\% of
the original length when the raw length ratio is at or below about $1.6\times$. Where it reaches
$2.17\times$, the scaling clamp binds on 59\,\% of segments and the result degrades. The cause was
diagnosed as unconstrained translation length rather than a fault in the retiming.

\textbf{Additional activity --- in-domain fine-tuning: completed, and correctly qualified.}
Fine-tuning raised BLEU by 2.34 and 3.23 with proper development-set checkpoint selection and a
single scored pass over the test split. The gain is register adaptation to the corpus's
conversational style, not improved translation adequacy, and is reported as such.

\section{Answers to the Research Questions}
\label{sec:answers}

\begin{description}[leftmargin=3.2em, style=nextline]
\item[\textbf{RQ1}] \textbf{86--95\,\% of the achievable speaker identity survives.} Measured
      against the same speaker's real recording in the other language, the system retains 94.5\,\%
      of the ceiling into Spanish and 86.1\,\% into English. The ceiling itself is only 0.445 under
      an ECAPA cosine, which is the fact that makes the raw score interpretable.

\item[\textbf{RQ2}] \textbf{Humans transfer prosody strongly; the system transfers a fraction of
      it.} Within-speaker, human cross-language correspondence is 0.849 for pitch level, 0.867 for
      duration, 0.601 for speaking rate and 0.214 for pitch range. The system reaches 0.311/0.314,
      0.352/0.610, 0.372/0.346 and 0.035/0.040 respectively. Pooled correlations are much higher
      but are largely between-speaker variance carried by voice cloning, and must not be quoted
      alone.

\item[\textbf{RQ3}] \textbf{Only partly.} Target-language pitch range is predictable from
      source-language prosody beyond a copy-plus-offset baseline, in both directions and under
      family-wise correction, but the effect is small (3.9--5.1\,\% of mean absolute error). Pitch
      level and duration ratio are not. For pitch level the regulariser selected
      $\lambda = 10^{6}$ and the model reduced itself to the baseline, which is the correct
      response to an absent signal.

\item[\textbf{RQ4}] \textbf{Recognition errors cost 1.84 and 3.34 BLEU}, in proportion to the
      recognition error rate in each direction. In-domain fine-tuning of the translation model
      recovers a comparable magnitude ($+2.34$ and $+3.23$ BLEU), but it does so by matching the
      corpus's register rather than by translating better, so the two effects should not be netted
      against each other.
\end{description}

\section{Contributions}
\label{sec:final_contributions}

\begin{enumerate}
  \item \textbf{A calibrated identity measurement.} Speaker similarity expressed as a percentage of
        a human same-speaker cross-language ceiling, rather than as an unanchored cosine.
  \item \textbf{A quantified prosody transfer gap}, reported within-speaker so that the identity
        result is not counted twice, with a demonstration that the pooled figure overstates
        transfer by roughly a factor of two for pitch level.
  \item \textbf{A negative result with a mechanism.} Cross-lingual prosody is predictable in pitch
        range and not in pitch level or timing, and the sign of the learned coefficient explains
        why.
  \item \textbf{A measured cost for the cascade}, 1.8--3.3 BLEU, obtained from the project's own
        data rather than cited.
  \item \textbf{A methodological correction} showing that correlating pitch in hertz without a
        plausibility gate understated F0 correspondence by roughly a factor of seven, together with
        evidence that the intuitive remedy --- narrowing the pitch tracker's search band --- makes
        the measurement worse.
  \item \textbf{A working, reproducible artefact}: a command-line research pipeline in which every
        reported number regenerates from committed result files, and a dubbing application that
        shares its model instances so that the two cannot drift apart.
\end{enumerate}

\section{Recommendations and Future Work}
\label{sec:future}

The following are ordered by the ratio of value to effort, as judged from the results above.

\begin{enumerate}
  \item \textbf{Conduct a listening study.} Twenty clips spanning good and poor cases, five to ten
        bilingual listeners, rating naturalness, speaker similarity and prosodic appropriateness on
        a five-point scale, with the automatic speaker-similarity scores correlated against the
        human similarity ratings. This validates the instrument as well as the system and closes
        the most significant gap in the present study.
  \item \textbf{Train a length-controlled translation model.} Prepend a length-class token to each
        source sentence during fine-tuning, derived from the target-to-source character ratio, and
        request a short rendition when the available slot is tight. This is the same twelve-minute
        fine-tune already performed, it attacks the $1.5\times$ duration inflation at its source
        rather than in post-processing, and it is measurable in four ways: BLEU,
        length-compliance rate, clamp rate, and final duration match.
  \item \textbf{Repair the two application defects identified in
        Section~\ref{sec:dubbing_result}}: truncate an overflowing segment at the next segment's
        start instead of summing overlapping audio, and fill stretches where no speech was detected
        with the ducked original audio instead of digital silence.
  \item \textbf{Compare F0 contours frame by frame under dynamic time warping.} Utterance-level
        mean and standard deviation cannot capture contour shape; a warped frame-level comparison
        would be a stronger instrument and would make these results directly comparable with
        published expressive-dubbing scores.
  \item \textbf{Close the loop on the prosody model.} The pitch-range relationship is real but
        currently only descriptive. Using the prediction to condition synthesis --- and measuring
        whether within-speaker correspondence improves --- would convert an analysis into an
        intervention.
  \item \textbf{Merge short recogniser segments before translation.} Sub-second slots produce
        impossible scaling demands in the dubbing application; merging adjacent short segments
        would remove a large share of the clamped cases at low cost.
  \item \textbf{Extend beyond one corpus and one language pair.} Replicating the ceiling
        methodology on a second same-speaker parallel corpus, or a typologically different language
        pair, would establish how far these figures generalise.
\end{enumerate}

\section{Closing Statement}
\label{sec:closing}

The most durable outcome of this project is not the pipeline, which is assembled from openly
available components, but the measurement it makes possible. A same-speaker parallel corpus turns
speech synthesis evaluation from a set of uncalibrated similarity scores into a comparison against
what the same person actually did. Under that comparison, the answer is specific: current zero-shot
cross-lingual synthesis carries the speaker, not the speech act. Identity survives at 86--95\,\% of
what is achievable; utterance-level prosody survives at roughly a third for pitch level and not at
all for pitch range. Stating that gap in reusable units is what this dissertation contributes, and
it is the quantity that the next generation of expressive dubbing systems will have to move.
